\begin{table}[!htbp]\centering\scriptsize
\caption{Descriptive incentive groups and anchor-matched compliance comparisons (panel 1 of 9)}
\begin{threeparttable}
\begin{tabular}{p{0.40\linewidth}rrrr}
\toprule
Group, period, scope, variant and statistic & N & Firms & Estimate & 95\% CI \\
\midrule
Only pre [0.60,0.80] All unique firms & 86 & 18 & 18.0000 & -- \\
Plus pre [0.60,0.80] All unique firms & 140 & 33 & 33.0000 & -- \\
Only post [0.60,0.80] All unique firms & 80 & 19 & 19.0000 & -- \\
Plus post [0.60,0.80] All unique firms & 119 & 32 & 32.0000 & -- \\
Only pre [0.60,0.80] All share 068 070 & 86 & 18 & 0.0465 & [0.0000, 0.1061] \\
Plus pre [0.60,0.80] All share 068 070 & 140 & 33 & 0.0643 & [0.0174, 0.1226] \\
Only post [0.60,0.80] All share 068 070 & 80 & 19 & 0.2875 & [0.1351, 0.4365] \\
Plus post [0.60,0.80] All share 068 070 & 119 & 32 & 0.1261 & [0.0444, 0.2097] \\
Only post minus pre [0.60,0.80] All share 068 070 & 80 / 86 & 19 / 18 & 0.2410 & [0.0946, 0.3805] \\
Plus post minus pre [0.60,0.80] All share 068 070 & 119 / 140 & 32 / 33 & 0.0618 & [-0.0260, 0.1581] \\
\bottomrule
\end{tabular}
\begin{tablenotes}[flushleft]\footnotesize
\item[] These are mechanical accounting counterfactuals, not attainable balance sheets, causal effects, or economic leverage. Minority interests may be genuine equity. Official RL1 is N/(N+E), N=liabilities minus advances, E=total equity including minority interests. N and E must be positive. Frozen MI replaces current minority interests with their end-2019 level, holding all other observed items at their realized values; the pre-period placebo freezes them at end-2016. Ex-MI is mechanical, not a deconsolidated balance sheet. Nonpositive total counterfactual denominators are unavailable. All shares are proportions. N and firms count available observations and unique firms for each estimate; contributing contrast counts are identified in the table notes. Both variants are reported: all observations and exclude jump dates (known half-year minority-interest jumps above 50 percent of current equity). Unknown flags are retained; later dates retain the jump in the MI stock. Percentile 95 percent intervals use 999 whole-firm resamples with shared draws, full histories and complete re-estimation. Intervals are pointwise; no significance stars. HK changes assume unchanged reporting currency within firm; currency metadata remain unresolved. Only denotes RL1-only breaches and Plus denotes RL1 plus other breaches at end-2019, requiring observed RL2 and RL3. Counts in [0.60,0.80] are listed first. D and fail given official use post 2020H2--2022H2 with the 2019 anchor and placebo 2017H2--2019H2 with the 2016 anchor. Share 068 070 uses [0.68,0.70) and pre 2017H1--2019H2. No pre-period D with a 2019 anchor is reported. For each statistic, scope, variant and period, both groups must contribute at least 15 unique firms with finite outcomes; contrasts require all contributing cells to pass. Otherwise inference is descriptive small cell and the CI is --. Contrasts list post / comparison counts; double differences are (Only post minus placebo) minus (Plus post minus placebo), with Only counts above Plus counts. Shared firm draws give paired percentile CIs. The economic value B of a relaxed cap need not be larger for Only firms; ranking it requires an assumption about financing constraints. Tier description: \citet{jin2025deleveraging}.
\end{tablenotes}
\end{threeparttable}
\end{table}

\begin{table}[!htbp]\centering\scriptsize
\caption{Descriptive incentive groups and anchor-matched compliance comparisons (panel 2 of 9)}
\begin{threeparttable}
\begin{tabular}{p{0.40\linewidth}rrrr}
\toprule
Group, period, scope, variant and statistic & N & Firms & Estimate & 95\% CI \\
\midrule
Only placebo [0.60,0.80] All D & 74 & 18 & 0.0405 & [0.0000, 0.1061] \\
Plus placebo [0.60,0.80] All D & 108 & 29 & -0.0185 & [-0.0630, 0.0000] \\
Only post [0.60,0.80] All D & 80 & 19 & 0.3125 & [0.1493, 0.4754] \\
Plus post [0.60,0.80] All D & 117 & 31 & 0.1538 & [0.0454, 0.2828] \\
Only post minus placebo [0.60,0.80] All D & 80 / 74 & 19 / 18 & 0.2720 & [0.0951, 0.4457] \\
Plus post minus placebo [0.60,0.80] All D & 117 / 108 & 31 / 29 & 0.1724 & [0.0548, 0.2967] \\
Only minus Plus post minus placebo [0.60,0.80] All D & \shortstack{80 / 74\\117 / 108} & \shortstack{19 / 18\\31 / 29} & 0.0996 & [-0.1117, 0.2980] \\
Only placebo [0.60,0.80] All fail given official & 12 & 5 & 0.2500 & -- \\
Plus placebo [0.60,0.80] All fail given official & 6 & 4 & 0.0000 & -- \\
Only post [0.60,0.80] All fail given official & 26 & 9 & 0.9615 & -- \\
\bottomrule
\end{tabular}
\begin{tablenotes}[flushleft]\footnotesize
\item[] These are mechanical accounting counterfactuals, not attainable balance sheets, causal effects, or economic leverage. Minority interests may be genuine equity. Official RL1 is N/(N+E), N=liabilities minus advances, E=total equity including minority interests. N and E must be positive. Frozen MI replaces current minority interests with their end-2019 level, holding all other observed items at their realized values; the pre-period placebo freezes them at end-2016. Ex-MI is mechanical, not a deconsolidated balance sheet. Nonpositive total counterfactual denominators are unavailable. All shares are proportions. N and firms count available observations and unique firms for each estimate; contributing contrast counts are identified in the table notes. Both variants are reported: all observations and exclude jump dates (known half-year minority-interest jumps above 50 percent of current equity). Unknown flags are retained; later dates retain the jump in the MI stock. Percentile 95 percent intervals use 999 whole-firm resamples with shared draws, full histories and complete re-estimation. Intervals are pointwise; no significance stars. HK changes assume unchanged reporting currency within firm; currency metadata remain unresolved. Only denotes RL1-only breaches and Plus denotes RL1 plus other breaches at end-2019, requiring observed RL2 and RL3. Counts in [0.60,0.80] are listed first. D and fail given official use post 2020H2--2022H2 with the 2019 anchor and placebo 2017H2--2019H2 with the 2016 anchor. Share 068 070 uses [0.68,0.70) and pre 2017H1--2019H2. No pre-period D with a 2019 anchor is reported. For each statistic, scope, variant and period, both groups must contribute at least 15 unique firms with finite outcomes; contrasts require all contributing cells to pass. Otherwise inference is descriptive small cell and the CI is --. Contrasts list post / comparison counts; double differences are (Only post minus placebo) minus (Plus post minus placebo), with Only counts above Plus counts. Shared firm draws give paired percentile CIs. The economic value B of a relaxed cap need not be larger for Only firms; ranking it requires an assumption about financing constraints. Tier description: \citet{jin2025deleveraging}.
\end{tablenotes}
\end{threeparttable}
\end{table}

\begin{table}[!htbp]\centering\scriptsize
\caption{Descriptive incentive groups and anchor-matched compliance comparisons (panel 3 of 9)}
\begin{threeparttable}
\begin{tabular}{p{0.40\linewidth}rrrr}
\toprule
Group, period, scope, variant and statistic & N & Firms & Estimate & 95\% CI \\
\midrule
Plus post [0.60,0.80] All fail given official & 28 & 10 & 0.6786 & -- \\
Only post minus placebo [0.60,0.80] All fail given official & 26 / 12 & 9 / 5 & 0.7115 & -- \\
Plus post minus placebo [0.60,0.80] All fail given official & 28 / 6 & 10 / 4 & 0.6786 & -- \\
Only minus Plus post minus placebo [0.60,0.80] All fail given official & \shortstack{26 / 12\\28 / 6} & \shortstack{9 / 5\\10 / 4} & 0.0330 & -- \\
Only pre all dates All share 068 070 & 119 & 20 & 0.0336 & [0.0000, 0.0774] \\
Plus pre all dates All share 068 070 & 255 & 45 & 0.0353 & [0.0084, 0.0738] \\
Only post all dates All share 068 070 & 95 & 20 & 0.2421 & [0.1091, 0.3751] \\
Plus post all dates All share 068 070 & 194 & 43 & 0.0773 & [0.0260, 0.1355] \\
Only post minus pre all dates All share 068 070 & 95 / 119 & 20 / 20 & 0.2085 & [0.0754, 0.3311] \\
Plus post minus pre all dates All share 068 070 & 194 / 255 & 43 / 45 & 0.0420 & [-0.0116, 0.1025] \\
\bottomrule
\end{tabular}
\begin{tablenotes}[flushleft]\footnotesize
\item[] These are mechanical accounting counterfactuals, not attainable balance sheets, causal effects, or economic leverage. Minority interests may be genuine equity. Official RL1 is N/(N+E), N=liabilities minus advances, E=total equity including minority interests. N and E must be positive. Frozen MI replaces current minority interests with their end-2019 level, holding all other observed items at their realized values; the pre-period placebo freezes them at end-2016. Ex-MI is mechanical, not a deconsolidated balance sheet. Nonpositive total counterfactual denominators are unavailable. All shares are proportions. N and firms count available observations and unique firms for each estimate; contributing contrast counts are identified in the table notes. Both variants are reported: all observations and exclude jump dates (known half-year minority-interest jumps above 50 percent of current equity). Unknown flags are retained; later dates retain the jump in the MI stock. Percentile 95 percent intervals use 999 whole-firm resamples with shared draws, full histories and complete re-estimation. Intervals are pointwise; no significance stars. HK changes assume unchanged reporting currency within firm; currency metadata remain unresolved. Only denotes RL1-only breaches and Plus denotes RL1 plus other breaches at end-2019, requiring observed RL2 and RL3. Counts in [0.60,0.80] are listed first. D and fail given official use post 2020H2--2022H2 with the 2019 anchor and placebo 2017H2--2019H2 with the 2016 anchor. Share 068 070 uses [0.68,0.70) and pre 2017H1--2019H2. No pre-period D with a 2019 anchor is reported. For each statistic, scope, variant and period, both groups must contribute at least 15 unique firms with finite outcomes; contrasts require all contributing cells to pass. Otherwise inference is descriptive small cell and the CI is --. Contrasts list post / comparison counts; double differences are (Only post minus placebo) minus (Plus post minus placebo), with Only counts above Plus counts. Shared firm draws give paired percentile CIs. The economic value B of a relaxed cap need not be larger for Only firms; ranking it requires an assumption about financing constraints. Tier description: \citet{jin2025deleveraging}.
\end{tablenotes}
\end{threeparttable}
\end{table}

\begin{table}[!htbp]\centering\scriptsize
\caption{Descriptive incentive groups and anchor-matched compliance comparisons (panel 4 of 9)}
\begin{threeparttable}
\begin{tabular}{p{0.40\linewidth}rrrr}
\toprule
Group, period, scope, variant and statistic & N & Firms & Estimate & 95\% CI \\
\midrule
Only placebo all dates All D & 100 & 20 & 0.0300 & [0.0000, 0.0818] \\
Plus placebo all dates All D & 187 & 38 & -0.0107 & [-0.0368, 0.0000] \\
Only post all dates All D & 95 & 20 & 0.2632 & [0.1176, 0.4066] \\
Plus post all dates All D & 192 & 42 & 0.0938 & [0.0259, 0.1840] \\
Only post minus placebo all dates All D & 95 / 100 & 20 / 20 & 0.2332 & [0.0808, 0.3831] \\
Plus post minus placebo all dates All D & 192 / 187 & 42 / 38 & 0.1044 & [0.0316, 0.1911] \\
Only minus Plus post minus placebo all dates All D & \shortstack{95 / 100\\192 / 187} & \shortstack{20 / 20\\42 / 38} & 0.1287 & [-0.0420, 0.2920] \\
Only placebo all dates All fail given official & 13 & 5 & 0.2308 & -- \\
Plus placebo all dates All fail given official & 10 & 5 & 0.0000 & -- \\
Only post all dates All fail given official & 26 & 9 & 0.9615 & -- \\
\bottomrule
\end{tabular}
\begin{tablenotes}[flushleft]\footnotesize
\item[] These are mechanical accounting counterfactuals, not attainable balance sheets, causal effects, or economic leverage. Minority interests may be genuine equity. Official RL1 is N/(N+E), N=liabilities minus advances, E=total equity including minority interests. N and E must be positive. Frozen MI replaces current minority interests with their end-2019 level, holding all other observed items at their realized values; the pre-period placebo freezes them at end-2016. Ex-MI is mechanical, not a deconsolidated balance sheet. Nonpositive total counterfactual denominators are unavailable. All shares are proportions. N and firms count available observations and unique firms for each estimate; contributing contrast counts are identified in the table notes. Both variants are reported: all observations and exclude jump dates (known half-year minority-interest jumps above 50 percent of current equity). Unknown flags are retained; later dates retain the jump in the MI stock. Percentile 95 percent intervals use 999 whole-firm resamples with shared draws, full histories and complete re-estimation. Intervals are pointwise; no significance stars. HK changes assume unchanged reporting currency within firm; currency metadata remain unresolved. Only denotes RL1-only breaches and Plus denotes RL1 plus other breaches at end-2019, requiring observed RL2 and RL3. Counts in [0.60,0.80] are listed first. D and fail given official use post 2020H2--2022H2 with the 2019 anchor and placebo 2017H2--2019H2 with the 2016 anchor. Share 068 070 uses [0.68,0.70) and pre 2017H1--2019H2. No pre-period D with a 2019 anchor is reported. For each statistic, scope, variant and period, both groups must contribute at least 15 unique firms with finite outcomes; contrasts require all contributing cells to pass. Otherwise inference is descriptive small cell and the CI is --. Contrasts list post / comparison counts; double differences are (Only post minus placebo) minus (Plus post minus placebo), with Only counts above Plus counts. Shared firm draws give paired percentile CIs. The economic value B of a relaxed cap need not be larger for Only firms; ranking it requires an assumption about financing constraints. Tier description: \citet{jin2025deleveraging}.
\end{tablenotes}
\end{threeparttable}
\end{table}

\begin{table}[!htbp]\centering\scriptsize
\caption{Descriptive incentive groups and anchor-matched compliance comparisons (panel 5 of 9)}
\begin{threeparttable}
\begin{tabular}{p{0.40\linewidth}rrrr}
\toprule
Group, period, scope, variant and statistic & N & Firms & Estimate & 95\% CI \\
\midrule
Plus post all dates All fail given official & 30 & 10 & 0.6333 & -- \\
Only post minus placebo all dates All fail given official & 26 / 13 & 9 / 5 & 0.7308 & -- \\
Plus post minus placebo all dates All fail given official & 30 / 10 & 10 / 5 & 0.6333 & -- \\
Only minus Plus post minus placebo all dates All fail given official & \shortstack{26 / 13\\30 / 10} & \shortstack{9 / 5\\10 / 5} & 0.0974 & -- \\
Only pre [0.60,0.80] Excl. dates share 068 070 & 86 & 18 & 0.0465 & [0.0000, 0.1061] \\
Plus pre [0.60,0.80] Excl. dates share 068 070 & 140 & 33 & 0.0643 & [0.0174, 0.1226] \\
Only post [0.60,0.80] Excl. dates share 068 070 & 80 & 19 & 0.2875 & [0.1351, 0.4365] \\
Plus post [0.60,0.80] Excl. dates share 068 070 & 119 & 32 & 0.1261 & [0.0444, 0.2097] \\
Only post minus pre [0.60,0.80] Excl. dates share 068 070 & 80 / 86 & 19 / 18 & 0.2410 & [0.0946, 0.3805] \\
Plus post minus pre [0.60,0.80] Excl. dates share 068 070 & 119 / 140 & 32 / 33 & 0.0618 & [-0.0260, 0.1581] \\
\bottomrule
\end{tabular}
\begin{tablenotes}[flushleft]\footnotesize
\item[] These are mechanical accounting counterfactuals, not attainable balance sheets, causal effects, or economic leverage. Minority interests may be genuine equity. Official RL1 is N/(N+E), N=liabilities minus advances, E=total equity including minority interests. N and E must be positive. Frozen MI replaces current minority interests with their end-2019 level, holding all other observed items at their realized values; the pre-period placebo freezes them at end-2016. Ex-MI is mechanical, not a deconsolidated balance sheet. Nonpositive total counterfactual denominators are unavailable. All shares are proportions. N and firms count available observations and unique firms for each estimate; contributing contrast counts are identified in the table notes. Both variants are reported: all observations and exclude jump dates (known half-year minority-interest jumps above 50 percent of current equity). Unknown flags are retained; later dates retain the jump in the MI stock. Percentile 95 percent intervals use 999 whole-firm resamples with shared draws, full histories and complete re-estimation. Intervals are pointwise; no significance stars. HK changes assume unchanged reporting currency within firm; currency metadata remain unresolved. Only denotes RL1-only breaches and Plus denotes RL1 plus other breaches at end-2019, requiring observed RL2 and RL3. Counts in [0.60,0.80] are listed first. D and fail given official use post 2020H2--2022H2 with the 2019 anchor and placebo 2017H2--2019H2 with the 2016 anchor. Share 068 070 uses [0.68,0.70) and pre 2017H1--2019H2. No pre-period D with a 2019 anchor is reported. For each statistic, scope, variant and period, both groups must contribute at least 15 unique firms with finite outcomes; contrasts require all contributing cells to pass. Otherwise inference is descriptive small cell and the CI is --. Contrasts list post / comparison counts; double differences are (Only post minus placebo) minus (Plus post minus placebo), with Only counts above Plus counts. Shared firm draws give paired percentile CIs. The economic value B of a relaxed cap need not be larger for Only firms; ranking it requires an assumption about financing constraints. Tier description: \citet{jin2025deleveraging}.
\end{tablenotes}
\end{threeparttable}
\end{table}

\begin{table}[!htbp]\centering\scriptsize
\caption{Descriptive incentive groups and anchor-matched compliance comparisons (panel 6 of 9)}
\begin{threeparttable}
\begin{tabular}{p{0.40\linewidth}rrrr}
\toprule
Group, period, scope, variant and statistic & N & Firms & Estimate & 95\% CI \\
\midrule
Only placebo [0.60,0.80] Excl. dates D & 74 & 18 & 0.0405 & [0.0000, 0.1061] \\
Plus placebo [0.60,0.80] Excl. dates D & 108 & 29 & -0.0185 & [-0.0630, 0.0000] \\
Only post [0.60,0.80] Excl. dates D & 80 & 19 & 0.3125 & [0.1493, 0.4754] \\
Plus post [0.60,0.80] Excl. dates D & 117 & 31 & 0.1538 & [0.0454, 0.2828] \\
Only post minus placebo [0.60,0.80] Excl. dates D & 80 / 74 & 19 / 18 & 0.2720 & [0.0951, 0.4457] \\
Plus post minus placebo [0.60,0.80] Excl. dates D & 117 / 108 & 31 / 29 & 0.1724 & [0.0548, 0.2967] \\
Only minus Plus post minus placebo [0.60,0.80] Excl. dates D & \shortstack{80 / 74\\117 / 108} & \shortstack{19 / 18\\31 / 29} & 0.0996 & [-0.1117, 0.2980] \\
Only placebo [0.60,0.80] Excl. dates fail given official & 12 & 5 & 0.2500 & -- \\
Plus placebo [0.60,0.80] Excl. dates fail given official & 6 & 4 & 0.0000 & -- \\
Only post [0.60,0.80] Excl. dates fail given official & 26 & 9 & 0.9615 & -- \\
\bottomrule
\end{tabular}
\begin{tablenotes}[flushleft]\footnotesize
\item[] These are mechanical accounting counterfactuals, not attainable balance sheets, causal effects, or economic leverage. Minority interests may be genuine equity. Official RL1 is N/(N+E), N=liabilities minus advances, E=total equity including minority interests. N and E must be positive. Frozen MI replaces current minority interests with their end-2019 level, holding all other observed items at their realized values; the pre-period placebo freezes them at end-2016. Ex-MI is mechanical, not a deconsolidated balance sheet. Nonpositive total counterfactual denominators are unavailable. All shares are proportions. N and firms count available observations and unique firms for each estimate; contributing contrast counts are identified in the table notes. Both variants are reported: all observations and exclude jump dates (known half-year minority-interest jumps above 50 percent of current equity). Unknown flags are retained; later dates retain the jump in the MI stock. Percentile 95 percent intervals use 999 whole-firm resamples with shared draws, full histories and complete re-estimation. Intervals are pointwise; no significance stars. HK changes assume unchanged reporting currency within firm; currency metadata remain unresolved. Only denotes RL1-only breaches and Plus denotes RL1 plus other breaches at end-2019, requiring observed RL2 and RL3. Counts in [0.60,0.80] are listed first. D and fail given official use post 2020H2--2022H2 with the 2019 anchor and placebo 2017H2--2019H2 with the 2016 anchor. Share 068 070 uses [0.68,0.70) and pre 2017H1--2019H2. No pre-period D with a 2019 anchor is reported. For each statistic, scope, variant and period, both groups must contribute at least 15 unique firms with finite outcomes; contrasts require all contributing cells to pass. Otherwise inference is descriptive small cell and the CI is --. Contrasts list post / comparison counts; double differences are (Only post minus placebo) minus (Plus post minus placebo), with Only counts above Plus counts. Shared firm draws give paired percentile CIs. The economic value B of a relaxed cap need not be larger for Only firms; ranking it requires an assumption about financing constraints. Tier description: \citet{jin2025deleveraging}.
\end{tablenotes}
\end{threeparttable}
\end{table}

\begin{table}[!htbp]\centering\scriptsize
\caption{Descriptive incentive groups and anchor-matched compliance comparisons (panel 7 of 9)}
\begin{threeparttable}
\begin{tabular}{p{0.40\linewidth}rrrr}
\toprule
Group, period, scope, variant and statistic & N & Firms & Estimate & 95\% CI \\
\midrule
Plus post [0.60,0.80] Excl. dates fail given official & 28 & 10 & 0.6786 & -- \\
Only post minus placebo [0.60,0.80] Excl. dates fail given official & 26 / 12 & 9 / 5 & 0.7115 & -- \\
Plus post minus placebo [0.60,0.80] Excl. dates fail given official & 28 / 6 & 10 / 4 & 0.6786 & -- \\
Only minus Plus post minus placebo [0.60,0.80] Excl. dates fail given official & \shortstack{26 / 12\\28 / 6} & \shortstack{9 / 5\\10 / 4} & 0.0330 & -- \\
Only pre all dates Excl. dates share 068 070 & 119 & 20 & 0.0336 & [0.0000, 0.0774] \\
Plus pre all dates Excl. dates share 068 070 & 255 & 45 & 0.0353 & [0.0084, 0.0738] \\
Only post all dates Excl. dates share 068 070 & 93 & 20 & 0.2473 & [0.1129, 0.3824] \\
Plus post all dates Excl. dates share 068 070 & 190 & 43 & 0.0789 & [0.0270, 0.1381] \\
Only post minus pre all dates Excl. dates share 068 070 & 93 / 119 & 20 / 20 & 0.2137 & [0.0820, 0.3386] \\
Plus post minus pre all dates Excl. dates share 068 070 & 190 / 255 & 43 / 45 & 0.0437 & [-0.0108, 0.1042] \\
\bottomrule
\end{tabular}
\begin{tablenotes}[flushleft]\footnotesize
\item[] These are mechanical accounting counterfactuals, not attainable balance sheets, causal effects, or economic leverage. Minority interests may be genuine equity. Official RL1 is N/(N+E), N=liabilities minus advances, E=total equity including minority interests. N and E must be positive. Frozen MI replaces current minority interests with their end-2019 level, holding all other observed items at their realized values; the pre-period placebo freezes them at end-2016. Ex-MI is mechanical, not a deconsolidated balance sheet. Nonpositive total counterfactual denominators are unavailable. All shares are proportions. N and firms count available observations and unique firms for each estimate; contributing contrast counts are identified in the table notes. Both variants are reported: all observations and exclude jump dates (known half-year minority-interest jumps above 50 percent of current equity). Unknown flags are retained; later dates retain the jump in the MI stock. Percentile 95 percent intervals use 999 whole-firm resamples with shared draws, full histories and complete re-estimation. Intervals are pointwise; no significance stars. HK changes assume unchanged reporting currency within firm; currency metadata remain unresolved. Only denotes RL1-only breaches and Plus denotes RL1 plus other breaches at end-2019, requiring observed RL2 and RL3. Counts in [0.60,0.80] are listed first. D and fail given official use post 2020H2--2022H2 with the 2019 anchor and placebo 2017H2--2019H2 with the 2016 anchor. Share 068 070 uses [0.68,0.70) and pre 2017H1--2019H2. No pre-period D with a 2019 anchor is reported. For each statistic, scope, variant and period, both groups must contribute at least 15 unique firms with finite outcomes; contrasts require all contributing cells to pass. Otherwise inference is descriptive small cell and the CI is --. Contrasts list post / comparison counts; double differences are (Only post minus placebo) minus (Plus post minus placebo), with Only counts above Plus counts. Shared firm draws give paired percentile CIs. The economic value B of a relaxed cap need not be larger for Only firms; ranking it requires an assumption about financing constraints. Tier description: \citet{jin2025deleveraging}.
\end{tablenotes}
\end{threeparttable}
\end{table}

\begin{table}[!htbp]\centering\scriptsize
\caption{Descriptive incentive groups and anchor-matched compliance comparisons (panel 8 of 9)}
\begin{threeparttable}
\begin{tabular}{p{0.40\linewidth}rrrr}
\toprule
Group, period, scope, variant and statistic & N & Firms & Estimate & 95\% CI \\
\midrule
Only placebo all dates Excl. dates D & 100 & 20 & 0.0300 & [0.0000, 0.0818] \\
Plus placebo all dates Excl. dates D & 187 & 38 & -0.0107 & [-0.0368, 0.0000] \\
Only post all dates Excl. dates D & 93 & 20 & 0.2688 & [0.1193, 0.4168] \\
Plus post all dates Excl. dates D & 188 & 42 & 0.0957 & [0.0266, 0.1872] \\
Only post minus placebo all dates Excl. dates D & 93 / 100 & 20 / 20 & 0.2388 & [0.0835, 0.3918] \\
Plus post minus placebo all dates Excl. dates D & 188 / 187 & 42 / 38 & 0.1064 & [0.0324, 0.1943] \\
Only minus Plus post minus placebo all dates Excl. dates D & \shortstack{93 / 100\\188 / 187} & \shortstack{20 / 20\\42 / 38} & 0.1324 & [-0.0421, 0.3019] \\
Only placebo all dates Excl. dates fail given official & 13 & 5 & 0.2308 & -- \\
Plus placebo all dates Excl. dates fail given official & 10 & 5 & 0.0000 & -- \\
Only post all dates Excl. dates fail given official & 26 & 9 & 0.9615 & -- \\
\bottomrule
\end{tabular}
\begin{tablenotes}[flushleft]\footnotesize
\item[] These are mechanical accounting counterfactuals, not attainable balance sheets, causal effects, or economic leverage. Minority interests may be genuine equity. Official RL1 is N/(N+E), N=liabilities minus advances, E=total equity including minority interests. N and E must be positive. Frozen MI replaces current minority interests with their end-2019 level, holding all other observed items at their realized values; the pre-period placebo freezes them at end-2016. Ex-MI is mechanical, not a deconsolidated balance sheet. Nonpositive total counterfactual denominators are unavailable. All shares are proportions. N and firms count available observations and unique firms for each estimate; contributing contrast counts are identified in the table notes. Both variants are reported: all observations and exclude jump dates (known half-year minority-interest jumps above 50 percent of current equity). Unknown flags are retained; later dates retain the jump in the MI stock. Percentile 95 percent intervals use 999 whole-firm resamples with shared draws, full histories and complete re-estimation. Intervals are pointwise; no significance stars. HK changes assume unchanged reporting currency within firm; currency metadata remain unresolved. Only denotes RL1-only breaches and Plus denotes RL1 plus other breaches at end-2019, requiring observed RL2 and RL3. Counts in [0.60,0.80] are listed first. D and fail given official use post 2020H2--2022H2 with the 2019 anchor and placebo 2017H2--2019H2 with the 2016 anchor. Share 068 070 uses [0.68,0.70) and pre 2017H1--2019H2. No pre-period D with a 2019 anchor is reported. For each statistic, scope, variant and period, both groups must contribute at least 15 unique firms with finite outcomes; contrasts require all contributing cells to pass. Otherwise inference is descriptive small cell and the CI is --. Contrasts list post / comparison counts; double differences are (Only post minus placebo) minus (Plus post minus placebo), with Only counts above Plus counts. Shared firm draws give paired percentile CIs. The economic value B of a relaxed cap need not be larger for Only firms; ranking it requires an assumption about financing constraints. Tier description: \citet{jin2025deleveraging}.
\end{tablenotes}
\end{threeparttable}
\end{table}

\begin{table}[!htbp]\centering\scriptsize
\caption{Descriptive incentive groups and anchor-matched compliance comparisons (panel 9 of 9)}
\begin{threeparttable}
\begin{tabular}{p{0.40\linewidth}rrrr}
\toprule
Group, period, scope, variant and statistic & N & Firms & Estimate & 95\% CI \\
\midrule
Plus post all dates Excl. dates fail given official & 30 & 10 & 0.6333 & -- \\
Only post minus placebo all dates Excl. dates fail given official & 26 / 13 & 9 / 5 & 0.7308 & -- \\
Plus post minus placebo all dates Excl. dates fail given official & 30 / 10 & 10 / 5 & 0.6333 & -- \\
Only minus Plus post minus placebo all dates Excl. dates fail given official & \shortstack{26 / 13\\30 / 10} & \shortstack{9 / 5\\10 / 5} & 0.0974 & -- \\
\bottomrule
\end{tabular}
\begin{tablenotes}[flushleft]\footnotesize
\item[] These are mechanical accounting counterfactuals, not attainable balance sheets, causal effects, or economic leverage. Minority interests may be genuine equity. Official RL1 is N/(N+E), N=liabilities minus advances, E=total equity including minority interests. N and E must be positive. Frozen MI replaces current minority interests with their end-2019 level, holding all other observed items at their realized values; the pre-period placebo freezes them at end-2016. Ex-MI is mechanical, not a deconsolidated balance sheet. Nonpositive total counterfactual denominators are unavailable. All shares are proportions. N and firms count available observations and unique firms for each estimate; contributing contrast counts are identified in the table notes. Both variants are reported: all observations and exclude jump dates (known half-year minority-interest jumps above 50 percent of current equity). Unknown flags are retained; later dates retain the jump in the MI stock. Percentile 95 percent intervals use 999 whole-firm resamples with shared draws, full histories and complete re-estimation. Intervals are pointwise; no significance stars. HK changes assume unchanged reporting currency within firm; currency metadata remain unresolved. Only denotes RL1-only breaches and Plus denotes RL1 plus other breaches at end-2019, requiring observed RL2 and RL3. Counts in [0.60,0.80] are listed first. D and fail given official use post 2020H2--2022H2 with the 2019 anchor and placebo 2017H2--2019H2 with the 2016 anchor. Share 068 070 uses [0.68,0.70) and pre 2017H1--2019H2. No pre-period D with a 2019 anchor is reported. For each statistic, scope, variant and period, both groups must contribute at least 15 unique firms with finite outcomes; contrasts require all contributing cells to pass. Otherwise inference is descriptive small cell and the CI is --. Contrasts list post / comparison counts; double differences are (Only post minus placebo) minus (Plus post minus placebo), with Only counts above Plus counts. Shared firm draws give paired percentile CIs. The economic value B of a relaxed cap need not be larger for Only firms; ranking it requires an assumption about financing constraints. Tier description: \citet{jin2025deleveraging}.
\end{tablenotes}
\end{threeparttable}
\end{table}

